\documentclass[conference]{IEEEtran}
\IEEEoverridecommandlockouts
\usepackage{cite}
\usepackage{amsmath,amssymb,amsfonts}
\usepackage{algorithmic}
\usepackage{graphicx}
\usepackage{textcomp}
\usepackage{xcolor}
\usepackage{booktabs}
\usepackage{multirow}
\usepackage{hyperref}
\usepackage{url}
\usepackage{caption}
\usepackage{subcaption}

\def\BibTeX{{\rm B\kern-.05em{\sc i\kern-.025em b}\kern-.08em
    T\kern-.1667em\lower.7ex\hbox{E}\kern-.125emX}}

\begin{document}

\title{Decentralized Edge AI Environmental Hazard Network: \\
First Prototype Stage (PS) Architecture and Feasibility}

\author{
\IEEEauthorblockN{Abhinav P. Nair}
\IEEEauthorblockA{
\textit{Department of Electrical Engineering} \\
\textit{Indian Institute of Technology Mandi} \\
Mandi, Himachal Pradesh, India \\
}
}

\maketitle

\begin{abstract}
Traditional meteorological monitoring relies on centralized, high-cost
macro-stations that report generalized weather data over broad geographic
radii. This paradigm degrades sharply in mountainous topographies, where
elevation change, ridge shadowing, and valley confinement create sharply
isolated microclimates. We present the First Prototype Stage (PS) of a
Decentralized Edge AI Environmental Hazard Network: a single, standalone,
solar-powered Edge AI node built around the ESP32-S3, a 12-parameter
environmental and hazard sensor suite, an on-device boosted-tree inference
pipeline, and a LoRa (865 MHz) multi-hop mesh radio layer. The node performs
sensor fusion and anomaly classification locally, transmitting only
event-driven, low-bandwidth alerts rather than raw telemetry, and remains
fully operational during cellular and internet blackouts. We detail the
compute, sensing, intelligence, and communication layers of the
architecture, contrast the design against commercial monitoring stations,
and present a bill-of-materials-based cost analysis showing per-node cost
of approximately INR 19,055 (full hazard configuration) or INR 6,600
(baseline configuration), against INR 1,20,000--2,50,000 for comparable
commercial units. We conclude with a scaling and optimization roadmap
toward a custom 4-layer PCB and mesh-wide deployment.
\end{abstract}

\begin{IEEEkeywords}
Edge AI, LoRa mesh network, environmental sensing, microclimate monitoring,
disaster early-warning, ESP32-S3, TinyML, wildfire detection, flash flood
prediction
\end{IEEEkeywords}

\section{Introduction}
\label{sec:intro}

Traditional meteorological monitoring relies on centralized, high-cost
macro-stations that provide generalized weather data for a broad geographic
radius. While effective for flat terrain, this paradigm fundamentally fails
in mountainous topographies characterized by complex elevation changes,
steep ridges, and deep valleys. The unique topography surrounding the IIT
Mandi campus and the greater Himachal Pradesh region creates isolated
microclimates, where weather on one side of a ridge can drastically differ
from the other.

The First Prototype Stage (PS) of the Edge AI Portable Weather Station
project aims to validate a paradigm shift: moving from centralized data
collection to a decentralized, self-healing swarm of hyper-local
intelligence nodes. The objective of the PS is to design, fabricate, and
test a single, standalone Edge AI node that integrates high-fidelity
atmospheric and environmental hazard sensors.

By executing machine learning models directly on the microcontroller (Edge
AI), the node locally processes raw sensor telemetry, identifies patterns
indicative of imminent environmental hazards such as flash floods or
approaching wildfire smoke, and transmits low-bandwidth alerts via a LoRa
mesh network. The PS serves as the fundamental building block: once the
standalone node's hardware, firmware, and power budget are validated, the
architecture can scale into a spatial mesh network for real-time tracking
of localized weather phenomena without reliance on vulnerable cellular
infrastructure.

\begin{figure}[htbp]
\centering
\includegraphics[width=\columnwidth]{../figures/system_architecture.png}
\caption{High-level system architecture of a single Edge AI hazard node,
showing the Compute, Sensing, Intelligence, and Communication layers.}
\label{fig:architecture}
\end{figure}

\section{Approach and System Architecture}
\label{sec:architecture}

Developing the PS requires a multi-disciplinary approach synthesizing
custom PCB design, ultra-low-power embedded systems, environmental
physics, and machine learning. To avoid compounding errors, the
architecture is isolated into four distinct layers: Compute, Sensing,
Intelligence, and Communication, as shown in Fig.~\ref{fig:architecture}.

\subsection{The Compute Core: ESP32-S3}
\label{subsec:compute}

The processing engine of the node is the ESP32-S3-WROOM-1 module. Selected
for its integrated Processor Instruction Extensions (PIE), the S3 provides
hardware-accelerated 128-bit vector math, critical for executing neural
networks and decision trees efficiently on battery power. The S3 offers 45
programmable GPIO pins, allowing simultaneous operation of an SPI bus (for
LoRa), an I2C bus (for multiple environmental sensors), hardware UART (for
optical particulate counters), and external hardware interrupts, all
buffered by 8 MB of PSRAM to handle dynamic mesh routing tables.

\subsection{The Sensor Array}
\label{subsec:sensors}

The PS is equipped with a robust sensor suite monitoring 12 distinct
parameters, categorized into baseline weather and severe environmental
hazards (Table~\ref{tab:sensors}).

\begin{itemize}
    \item \textbf{Core Meteorological Baseline:} Temperature, humidity, and
    barometric pressure via the Bosch BME280 (I2C). Wind speed and
    direction via the SparkFun Weather Meter Kit. Precipitation via a
    piezoelectric acoustic sensor (DFRobot SEN0575), avoiding the leveling
    requirement of tipping-bucket designs. UV Index via the LTR390.
    \item \textbf{Hazard and Air Quality Array:} A PMS5003 laser-scattering
    sensor (UART) for PM2.5/PM10 particulate counting. A Sensirion SCD41
    photoacoustic NDIR sensor for true CO\textsubscript{2}. An SGP41 for
    VOCs and NOx. A MICS-6814 for multi-gas detection (CO, NO2, NH3).
    \item \textbf{Predictive Hardware:} An AS3935 lightning sensor analyzes
    RF emissions to estimate distance to approaching storm fronts up to 40
    km away, providing a leading indicator for the predictive models.
    \item \textbf{System Telemetry:} On-board 18650 battery voltage and
    internal enclosure temperature (DS18B20) for power-cycle regulation
    and water-ingress detection.
\end{itemize}

\begin{table}[htbp]
\caption{Sensor Array Summary}
\label{tab:sensors}
\centering
\begin{tabular}{@{}p{2.1cm}p{2.9cm}p{2.2cm}@{}}
\toprule
\textbf{Category} & \textbf{Component} & \textbf{Interface} \\
\midrule
Temp/Hum/Pressure & Bosch BME280 & I2C \\
Wind Speed/Dir & SparkFun Weather Meter Kit & Analog/Digital \\
Precipitation & DFRobot SEN0575 (piezo) & Analog \\
UV Index & LTR390 & I2C \\
PM2.5 / PM10 & Plantower PMS5003 & UART \\
CO\textsubscript{2} & Sensirion SCD41 & I2C \\
VOC / NOx & Sensirion SGP41 & I2C \\
Multi-gas (CO, NO2, NH3) & MICS-6814 & Analog \\
Lightning distance & AS3935 & I2C/SPI \\
Enclosure temp & DS18B20 & 1-Wire \\
Battery voltage & Onboard ADC divider & Analog \\
\bottomrule
\end{tabular}
\end{table}

\subsection{Edge AI Implementation}
\label{subsec:edgeai}

Rather than functioning as a passive data logger, the PS acts as an
intelligent agent. Using historical weather data, regression models based
on boosting algorithms (XGBoost or AdaBoost) are trained offline in
Python. These models use feature combinations, e.g., rapid barometric
drops coupled with rising PM2.5 levels and decreasing lightning distances,
to classify and predict localized anomalies.

Using conversion toolkits such as \texttt{micromlgen}, the trained model
is exported as a raw C++ header and integrated into the ESP32-S3 firmware.
The Ultra-Low-Power (ULP) coprocessor continuously polls the sensors in
the background; only when specific thresholds are breached does the main
dual-core processor wake, execute the boosted-tree inference, and
determine whether an alert should be broadcast.

\begin{figure}[htbp]
\centering
\includegraphics[width=\columnwidth]{../figures/edge_ai_flowchart.png}
\caption{Edge AI decision flow: from ULP sensor polling to threshold
breach, wake, inference, and mesh alert broadcast.}
\label{fig:flowchart}
\end{figure}

\subsection{Communication Layer: LoRa Mesh}
\label{subsec:comms}

To achieve a 10 km operational range in mountainous terrain without
cellular dependency, the PS uses the Semtech SX1276 (RFM95W) transceiver
tuned to the legal 865 MHz ISM band. Rather than a direct point-to-point
link, which would be blocked by topology, the network employs a multi-hop
mesh routing protocol (e.g., Meshtastic or RadioHead). The node transmits
event-driven insights, using federated-learning principles to distribute
small AI parameter updates or threat alerts to neighboring nodes,
eventually routing data to a local base-station dashboard.

\subsection{Enclosure and Power Management}
\label{subsec:power}

Physical survival of the electronics is provided by an FDM 3D-printed
Stevenson screen in PETG filament, ensuring natural airflow while
shielding sensors from direct insolation and precipitation. Power is
sustained via a high-capacity 18650 lithium-ion cell, managed by a
dedicated TP4056 or CN3065 charge controller, and replenished by a 5 W
external solar panel.

\section{Difference from Commercial Monitoring Stations}
\label{sec:comparison}

The Edge AI mesh architecture diverges fundamentally from legacy
commercial systems, offering distinct operational advantages in
constrained, dynamic environments (Table~\ref{tab:comparison}).

\begin{table*}[htbp]
\caption{Architectural Comparison Between Legacy and Edge AI Weather Systems}
\label{tab:comparison}
\centering
\begin{tabular}{@{}p{2.6cm}p{6.8cm}p{6.8cm}@{}}
\toprule
\textbf{Feature Domain} & \textbf{Traditional Commercial Station} & \textbf{Edge AI Mesh Node} \\
\midrule
Data Resolution & Macro-level; generalized data over $>$50 km radius. & Hyper-local; maps precise spatial microclimates and topographical variation. \\
Processing Paradigm & Passive telemetry; raw, continuous streams to a central cloud server. & Edge inference; ingests raw data locally, runs ML models (e.g., XGBoost), transmits insights. \\
Connectivity & Fragile; relies on Wi-Fi, 4G, or Satellite IoT; fails during blackouts. & Zero-cloud survivability; ad-hoc LoRa mesh reroutes dynamically around failed nodes. \\
Bandwidth Usage & High; raw data matrices choke low-bandwidth networks. & Event-driven; small payload alerts only (e.g., ``Flood Risk 85\%''). \\
Deployment Logistics & Static, heavy; requires professional installation and leveling. & Portable, agile, solar-powered; rapid deployment along supply chains or ridges. \\
\bottomrule
\end{tabular}
\end{table*}

\subsection{Zero-Cloud Survivability}

The most significant differentiator is resilience. Commercial IoT
stations become inert when cellular networks fail during severe storms.
Because the PS runs its predictive algorithms natively on the ESP32-S3
bare metal, it retains full analytical capability offline. The LoRa mesh
acts as an independent intranet, bypassing ISP infrastructure to deliver
warnings directly to localized physical sirens or base stations.

\subsection{High-Fidelity Hazard Detection}

Commercial stations prioritize broad meteorological trends. By
integrating the PMS5003 laser particulate counter and Sensirion gas
arrays alongside acoustic precipitation and lightning RF detection, the
PS doubles as an environmental hazard node, tuned for wildfire smoke
tracking and localized flash-flood prediction.

\section{Cost Breakdown and Economic Viability}
\label{sec:cost}

Commercial weather stations with equivalent granular hazard detection
(PM2.5, true CO\textsubscript{2}, lightning distance, multi-gas)
typically cost between INR 1,20,000 and INR 2,50,000 per unit, rendering
dense spatial microclimate mapping financially impractical for
independent researchers or local municipalities. By using off-the-shelf
breakout modules and the open-source C++ ecosystem, the Edge AI PS
achieves comparable analytical density at a fraction of the cost
(Table~\ref{tab:bom}).

\begin{table}[htbp]
\caption{Estimated Bill of Materials for the PS Node}
\label{tab:bom}
\centering
\begin{tabular}{@{}p{2.3cm}p{3.4cm}r@{}}
\toprule
\textbf{Category} & \textbf{Components} & \textbf{Cost (INR)} \\
\midrule
Compute \& Radio & ESP32-S3 WROOM-1, RFM95W SX1276, passives & 1,430 \\
Core Meteorological & BME280, Weather Meter Kit, DS18B20, SEN0575, LTR390 & 3,350 \\
Particulate \& Gases & PMS5003, SGP41, MICS-6814, SCD41 & 9,055 \\
Specialized Sensing & AS3935 & 935 \\
Power Architecture & 18650 cell, 5W solar panel, CN3065 & 2,380 \\
Fabrication & PCB (5x batch), PETG, conformal coating & 1,905 \\
\midrule
\textbf{Total per Node} & \textbf{Complete research-grade node} & \textbf{19,055} \\
\bottomrule
\end{tabular}
\end{table}

\subsection{Scalability and Optimization}

At approximately INR 19,055 per node, a 20-node mesh network represents a
capital expenditure of INR 3.8 Lakhs, still substantially lower than a
two-unit commercial deployment. Transitioning from prototype to
production mesh, a unified schematic in KiCad and a custom 4-layer PCB
integrating all sensor ICs directly (bypassing redundant regulators and
LEDs found on commercial breakout boards) is projected to reduce per-node
cost by an additional 20--30\%. For baseline nodes where hazard gases are
not required, removing the SCD41, SGP41, and MICS-6814 reduces cost to
approximately INR 6,600, enabling mass spatial deployment while reserving
fully equipped INR 19k nodes for strategic choke points in the valley.

\section{Conclusion}
\label{sec:conclusion}

The First Prototype Stage of the Edge AI Portable Weather Station
represents a feasible, economically disruptive approach to microclimate
monitoring. By pairing the machine-learning acceleration of the ESP32-S3
with the offline resilience of an 865 MHz LoRa mesh network, the system
overcomes the geographic and infrastructural limitations of centralized
weather tracking. The immediate next phase involves finalizing hardware
integration of the 12-sensor array, capturing baseline tabular data, and
deploying the C++-converted boosting algorithm to validate power
consumption during localized inference loops. Once validated, the node is
ready to replicate and scale across mountainous topography, establishing
the foundation of a true swarm-intelligence environmental network.

\section*{Acknowledgment}
The author thanks the Prayas 4.0 program at IIT Mandi for lab access and
support.

\begin{thebibliography}{00}
\bibitem{b1} Bosch Sensortec, ``BME280 Combined Humidity and Pressure Sensor Datasheet,'' 2020.
\bibitem{b2} Semtech Corporation, ``SX1276/77/78/79 Datasheet,'' 2020.
\bibitem{b3} Espressif Systems, ``ESP32-S3 Series Datasheet,'' 2023.
\bibitem{b4} Sensirion AG, ``SCD41 CO2 Sensor Datasheet,'' 2021.
\bibitem{b5} T. Chen and C. Guestrin, ``XGBoost: A Scalable Tree Boosting System,'' \textit{Proc. 22nd ACM SIGKDD}, 2016.
\bibitem{b6} Meshtastic Project, ``Meshtastic Documentation,'' \url{https://meshtastic.org}, 2026.
\end{thebibliography}

\end{document}
